% !PRINTED 83
\weberpage{83}
%<<CONT
\weq{15}{\vartheta_{00}^2\,\vartheta_{00}^2(v)\,\frac{d^2\log\vartheta_{00}(v)}{dv^2} = \vartheta_{00}\vartheta_{00}''\vartheta_{00}^2(v) + \vartheta_{11}'^2\vartheta_{11}^2(v),}
\weq{16}{\vartheta_{01}^2\,\vartheta_{01}^2(v)\,\frac{d^2\log\vartheta_{01}(v)}{dv^2} = \vartheta_{01}\vartheta_{01}''\vartheta_{01}^2(v) - \vartheta_{11}'^2\vartheta_{11}^2(v),}
\weq{17}{\vartheta_{10}^2\,\vartheta_{10}^2(v)\,\frac{d^2\log\vartheta_{10}(v)}{dv^2} = \vartheta_{10}\vartheta_{10}''\vartheta_{10}^2(v) - \vartheta_{11}'^2\vartheta_{11}^2(v),}
\weq{18}{\vartheta_{00}^2\,\vartheta_{11}^2(v)\,\frac{d^2\log\vartheta_{11}(v)}{dv^2} = \vartheta_{00}\vartheta_{00}''\vartheta_{11}^2(v) - \vartheta_{11}'^2\vartheta_{00}^2(v).}

Hieraus lassen sich noch weitere Relationen herleiten durch fortgesetzte Differentiation. Wir führen noch eine dieser Formeln an, die sich ergibt, wenn man (15) noch zweimal nach $v$ differentiiert und dann $v=0$ setzt. Drückt man die Differentiationen nach $v$ durch solche nach $\omega$ aus mittels der partiellen Differentialgleichung (4), so folgt:
\[
\frac{d^2\log\vartheta_{00}}{d\omega^2} - 2\left(\frac{d\log\vartheta_{00}}{d\omega}\right)^2 = -\frac{\pi^2}{8}\,\vartheta_{10}^4\vartheta_{01}^4,
\]
oder
\weq{19}{\frac{d}{d\omega}\left(\frac{1}{\vartheta_{00}^4}\,\frac{d\vartheta_{00}^2}{d\omega}\right) = -\frac{\pi^2}{4}\,\frac{\vartheta_{10}^4\vartheta_{01}^4}{\vartheta_{00}^2}.}

\begin{center}
\textbf{\S\,24. \gesperrt{Darstellung der $\vartheta$-Funktionen durch unendliche}}\\
\textbf{\gesperrt{Produkte.}}
\end{center}

Die Theorie der $\vartheta$-Funktionen ist nun so weit gefördert, daß sich ihre Darstellung sehr leicht ergibt. Damit wird dann die Existenz dieser Funktionen nachgewiesen und die bisherigen Betrachtungen erhalten erst ihren sicheren Boden. Zwei Wege zu diesem Ziele stehen uns offen. Der erste geht aus von den uns schon bekannten Nullpunkten der $\vartheta$-Funktionen und setzt daraus unendliche Produkte zusammen.

Die Funktion $\vartheta_{00}(v)$ verschwindet nach \S\,21, wenn
\[
2v = (2\nu-1)\omega + (2\mu-1)
\]
ist, worin $\nu, \mu$ ganze Zahlen sind, oder wenn
\[
e^{\pm 2\pi i v} = -e^{\pi i \omega(2\nu-1)},
\]
worin wir nun $\nu$ auf \gesperrt{positive} Zahlen beschränken können. Setzen wir also zur Abkürzung
\[
e^{\pi i \omega} = q,
\]
%CONT>>
